\begin{flushleft}
\begin{tikzpicture}[every node/.style={font=\footnotesize},
  chSixRelationMatrix/.style={matrix of nodes,nodes in empty cells,
    nodes={draw=black,line width=.6pt,align=left,inner xsep=3pt,inner ysep=2pt,minimum height=7mm,anchor=center,text height=2.7mm,text depth=.8mm},
    column sep=-.6pt,row sep=-.6pt,row 1/.style={nodes={fill=chSixTint}}}]
\matrix (s) [chSixRelationMatrix,column 1/.style={nodes={text width=25mm}},column 2/.style={nodes={text width=19mm}}] {
StudentID (PK) & Name\\
101 & Rupa\\
102 & Nabil\\
103 & Tania\\
};
\matrix (p) [chSixRelationMatrix,anchor=north west,column 1/.style={nodes={text width=33mm}},column 2/.style={nodes={text width=17mm}}]
at ([xshift=20mm]s.north east) {
StudentID (PK, FK) & SSCGPA\\
101 & 5.00\\
102 & 4.75\\
103 & 4.50\\
};
\node[anchor=south west,font=\small] at ([yshift=2mm]s.north west) {\textbf{Student: 1}};
\node[anchor=south west,font=\small] at ([yshift=2mm]p.north west) {\textbf{Profile: 1}};
\foreach \r in {2,3,4}{\draw[bookarrow] (s-\r-2.east)--(p-\r-1.west);}
\end{tikzpicture}
\captionof{table}{একই StudentID-তে একটি Student ও সর্বোচ্চ একটি Profile-এর মিল}
\end{flushleft}
